% !TeX root = ../exemple-A-assignment.tex
% !TeX encoding = UTF-8
%% =====================================================================
%%  Exemple A — Devoir à la maison de physique théorique
%%  Sujet : l'oscillateur harmonique quantique à une dimension.
%%  Montre : rappels (définition, théorème, lemme, exemple, remarque),
%%  exercice à plusieurs questions, solutions détaillées, démonstrations,
%%  équations numérotées et référencées, notations de Dirac, unités SI.
%% =====================================================================

\begin{resume}
  Ce devoir étudie l'oscillateur harmonique quantique unidimensionnel par la
  méthode algébrique des opérateurs d'échelle. On établit le spectre de
  l'énergie, on calcule les dispersions en position et en impulsion, et on
  vérifie l'inégalité de Heisenberg–Robertson.
\end{resume}

%% ---------------------------------------------------------------------
\section{Rappels théoriques}
%% ---------------------------------------------------------------------

\subsection{Hamiltonien et opérateurs d'échelle}

Une particule de masse $m$ soumise au potentiel $V(x)=\tfrac12 m\omega^{2}x^{2}$
est décrite par le hamiltonien
\begin{equation}
  \op{H} = \frac{\op{p}^{2}}{2m} + \frac{1}{2}\,m\omega^{2}\,\op{x}^{2},
  \qquad \comm{\op{x}}{\op{p}} = \ii\hbar\,.
  \label{eq:H}
\end{equation}

\begin{definition}[Opérateurs d'échelle]\label{def:echelle}
  On appelle opérateur d'annihilation $\op{a}$ et opérateur de création
  $\op{a}\adj$ les opérateurs
  \begin{align}
    \op{a}  &= \sqrt{\frac{m\omega}{2\hbar}}\,\Bigl(\op{x} + \frac{\ii}{m\omega}\,\op{p}\Bigr),
    \label{eq:a}\\
    \op{a}\adj &= \sqrt{\frac{m\omega}{2\hbar}}\,\Bigl(\op{x} - \frac{\ii}{m\omega}\,\op{p}\Bigr),
    \label{eq:adag}
  \end{align}
  et l'on pose $\op{N}=\op{a}\adj\op{a}$ (opérateur nombre).
\end{definition}

\begin{lemme}\label{lem:comm}
  Les opérateurs d'échelle vérifient
  \begin{equation}
    \comm{\op{a}}{\op{a}\adj} = 1,
    \qquad
    \comm{\op{N}}{\op{a}} = -\op{a},
    \qquad
    \comm{\op{N}}{\op{a}\adj} = \op{a}\adj .
    \label{eq:comm}
  \end{equation}
\end{lemme}

\begin{demonstration}
  Comme $\comm{\op{x}}{\op{x}}=\comm{\op{p}}{\op{p}}=0$ et
  $\comm{\op{p}}{\op{x}}=-\ii\hbar$, la bilinéarité du commutateur donne
  \begin{align*}
    \comm{\op{a}}{\op{a}\adj}
      &= \frac{m\omega}{2\hbar}\Bigl(
         -\frac{\ii}{m\omega}\comm{\op{x}}{\op{p}}
         +\frac{\ii}{m\omega}\comm{\op{p}}{\op{x}}\Bigr)
       = \frac{1}{2\hbar}\bigl(-\ii\cdot\ii\hbar-\ii\cdot\ii\hbar\bigr)
       = 1 .
  \end{align*}
  Ensuite, l'identité $\comm{AB}{C}=A\comm{B}{C}+\comm{A}{C}B$ appliquée à
  $\op{N}=\op{a}\adj\op{a}$ fournit
  $\comm{\op{N}}{\op{a}}=\op{a}\adj\comm{\op{a}}{\op{a}}+\comm{\op{a}\adj}{\op{a}}\op{a}=-\op{a}$
  puisque $\comm{\op{a}\adj}{\op{a}}=-1$. Le dernier commutateur s'obtient de
  même : $\comm{\op{N}}{\op{a}\adj}=\op{a}\adj\comm{\op{a}}{\op{a}\adj}=\op{a}\adj$.
\end{demonstration}

\subsection{Relation d'incertitude}

Pour une observable $\op{A}$ et un état normalisé $\ket{\psi}$, on note
$\ev{\op{A}}{\psi}=\bra{\psi}\op{A}\ket{\psi}$ sa valeur moyenne et
$\Delta A=\sqrt{\ev{\op{A}^{2}}{\psi}-\ev{\op{A}}{\psi}^{2}}$ sa dispersion.

\begin{theoreme}[Inégalité de Robertson]\label{thm:robertson}
  Pour deux observables $\op{A}$ et $\op{B}$ et tout état normalisé
  $\ket{\psi}$,
  \begin{equation}
    \Delta A\,\Delta B \;\geq\; \frac{1}{2}\,\abs[\big]{\ev{\comm{\op{A}}{\op{B}}}{\psi}} .
    \label{eq:robertson}
  \end{equation}
\end{theoreme}

\begin{demonstration}
  Posons $\ket{f}=(\op{A}-\ev{\op{A}}{\psi})\ket{\psi}$ et
  $\ket{g}=(\op{B}-\ev{\op{B}}{\psi})\ket{\psi}$, de sorte que
  $\braket{f}{f}=(\Delta A)^{2}$ et $\braket{g}{g}=(\Delta B)^{2}$.
  L'inégalité de Cauchy–Schwarz s'écrit
  \begin{equation}
    (\Delta A)^{2}(\Delta B)^{2} \geq \abs{\braket{f}{g}}^{2}
    \geq \bigl(\Imag\braket{f}{g}\bigr)^{2}.
    \label{eq:CS}
  \end{equation}
  Or $\braket{g}{f}=\braket{f}{g}\cc$, d'où
  \begin{equation}
    \Imag\braket{f}{g}=\frac{\braket{f}{g}-\braket{g}{f}}{2\ii}
    =\frac{1}{2\ii}\,\ev{\comm{\op{A}}{\op{B}}}{\psi},
  \end{equation}
  car les termes en $\ev{\op{A}}{\psi}$ et $\ev{\op{B}}{\psi}$ se compensent
  dans la différence. Cette quantité étant réelle, son carré est le carré de
  son module, soit $\tfrac14\abs{\ev{\comm{\op{A}}{\op{B}}}{\psi}}^{2}$. En
  reportant dans l'\cref{eq:CS} et en prenant la racine carrée, on obtient
  l'\cref{eq:robertson}.
\end{demonstration}

\begin{exemple}[État fondamental]\label{ex:fondamental}
  L'état $\ket{0}$ défini par $\op{a}\ket{0}=0$ vérifie, en représentation
  position où $\op{p}=-\ii\hbar\,\dd/\dd x$,
  \begin{equation}
    x\,\psi_{0}(x)+\frac{\hbar}{m\omega}\dv{\psi_{0}}{x}=0
    \quad\Longrightarrow\quad
    \psi_{0}(x)=\Bigl(\frac{m\omega}{\pi\hbar}\Bigr)^{1/4}
      \exp\Bigl(-\frac{m\omega x^{2}}{2\hbar}\Bigr).
    \label{eq:psi0}
  \end{equation}
\end{exemple}

%% ---------------------------------------------------------------------
\section{Exercice : spectre et dispersions}
%% ---------------------------------------------------------------------

\begin{exercice}[Oscillateur harmonique quantique]\label{exo:oh}
  On reprend les notations de la \cref{def:echelle}. Les états propres
  normalisés de $\op{N}$ sont notés $\ket{n}$, avec $\op{N}\ket{n}=n\ket{n}$.
\end{exercice}

%% ----- Question 1 -----
\begin{question}[2 pts]\label{q:H}
  Exprimer $\op{H}$ en fonction de $\op{N}$ et de $\hbar\omega$.
\end{question}

\begin{solution}
  Calculons $\op{a}\adj\op{a}$ à partir des \cref{eq:a,eq:adag} :
  \begin{align}
    \op{a}\adj\op{a}
      &= \frac{m\omega}{2\hbar}\Bigl(\op{x}-\frac{\ii}{m\omega}\op{p}\Bigr)
         \Bigl(\op{x}+\frac{\ii}{m\omega}\op{p}\Bigr) \notag\\
      &= \frac{m\omega}{2\hbar}\Bigl(\op{x}^{2}+\frac{\op{p}^{2}}{m^{2}\omega^{2}}
         +\frac{\ii}{m\omega}\comm{\op{x}}{\op{p}}\Bigr)
       = \frac{\op{H}}{\hbar\omega}-\frac12,
    \label{eq:NH}
  \end{align}
  où l'on a utilisé $\comm{\op{x}}{\op{p}}=\ii\hbar$ et reconnu $\op{H}$
  dans l'\cref{eq:H}. Il en résulte
  \begin{equation}
    \boxed{\;\op{H}=\hbar\omega\Bigl(\op{N}+\tfrac12\Bigr)\;}
    \label{eq:Hdiag}
  \end{equation}
  Le hamiltonien est donc diagonal dans la base des $\ket{n}$.
\end{solution}

%% ----- Question 2 -----
\begin{question}[4 pts]\label{q:spectre}
  Montrer que le spectre de $\op{N}$ est l'ensemble $\NN$ des entiers naturels,
  et que
  \begin{equation}
    \op{a}\ket{n}=
    \begin{cases}
      \sqrt{n}\,\ket{n-1}, & n\geq 1,\\[2pt]
      0, & n=0,
    \end{cases}
    \qquad
    \op{a}\adj\ket{n}=\sqrt{n+1}\,\ket{n+1}.
    \label{eq:action}
  \end{equation}
  En déduire les niveaux d'énergie $E_{n}$.
\end{question}

\begin{solution}
  Soit $\ket{\nu}$ un état propre normalisé de $\op{N}$ de valeur propre
  $\nu$.

  \emph{(i) Positivité.} On a
  $\nu=\ev{\op{N}}{\nu}=\ev{\op{a}\adj\op{a}}{\nu}=\norm{\op{a}\ket{\nu}}^{2}\geq 0$,
  avec égalité si et seulement si $\op{a}\ket{\nu}=0$.

  \emph{(ii) Échelle.} D'après l'\cref{eq:comm},
  $\op{N}\op{a}\ket{\nu}=(\op{a}\op{N}-\op{a})\ket{\nu}=(\nu-1)\,\op{a}\ket{\nu}$ :
  si $\op{a}\ket{\nu}\neq0$, c'est un état propre de valeur propre $\nu-1$,
  dont le carré de la norme vaut $\nu$. De même
  $\op{N}\op{a}\adj\ket{\nu}=(\nu+1)\,\op{a}\adj\ket{\nu}$, avec
  $\norm{\op{a}\adj\ket{\nu}}^{2}=\ev{\op{a}\op{a}\adj}{\nu}=\nu+1>0$.

  \emph{(iii) Quantification.} Appliquons $\op{a}$ de façon répétée : les
  valeurs propres $\nu,\nu-1,\nu-2,\dots$ doivent rester positives ou
  nulles d'après (i). La suite s'interrompt seulement lorsqu'on atteint la
  valeur propre $0$ (où $\op{a}\ket{0}=0$), ce qui exige $\nu\in\NN$. Réciproquement,
  $\ket{n}\propto(\op{a}\adj)^{n}\ket{0}$ existe pour tout $n\in\NN$ d'après
  (ii). Les relations \eqref{eq:action} découlent des normes calculées en (ii),
  les phases étant choisies réelles positives.

  Par l'\cref{eq:Hdiag}, les niveaux d'énergie sont
  \begin{equation}
    E_{n}=\hbar\omega\Bigl(n+\frac12\Bigr),\qquad n\in\NN .
    \label{eq:En}
  \end{equation}
\end{solution}

%% ----- Question 3 -----
\begin{question}[4 pts]\label{q:dispersions}
  Calculer $\ev{\op{x}}{n}$, $\ev{\op{p}}{n}$, $\ev{\op{x}^{2}}{n}$ et
  $\ev{\op{p}^{2}}{n}$, puis le produit $\Delta x\,\Delta p$. Vérifier
  l'\cref{eq:robertson} pour $\op{A}=\op{x}$ et $\op{B}=\op{p}$.
\end{question}

\begin{solution}
  Inversons les \cref{eq:a,eq:adag} :
  \begin{equation}
    \op{x}=\sqrt{\frac{\hbar}{2m\omega}}\,\bigl(\op{a}+\op{a}\adj\bigr),
    \qquad
    \op{p}=-\ii\sqrt{\frac{m\hbar\omega}{2}}\,\bigl(\op{a}-\op{a}\adj\bigr).
    \label{eq:xp}
  \end{equation}
  Comme $\op{a}$ et $\op{a}\adj$ changent $n$ d'une unité, leurs valeurs
  moyennes dans $\ket{n}$ sont nulles : $\ev{\op{x}}{n}=\ev{\op{p}}{n}=0$.
  Pour les carrés, seuls les termes qui conservent $n$ survivent, et
  $\op{a}\op{a}\adj=\op{N}+1$ d'après l'\cref{eq:comm} :
  \begin{align}
    \ev{\op{x}^{2}}{n}
      &=\frac{\hbar}{2m\omega}\,\ev{\op{a}\op{a}\adj+\op{a}\adj\op{a}}{n}
       =\frac{\hbar}{2m\omega}\,(2n+1),
    \label{eq:x2}\\
    \ev{\op{p}^{2}}{n}
      &=\frac{m\hbar\omega}{2}\,\ev{\op{a}\op{a}\adj+\op{a}\adj\op{a}}{n}
       =\frac{m\hbar\omega}{2}\,(2n+1).
    \label{eq:p2}
  \end{align}
  Les dispersions valent donc $\Delta x=\sqrt{\ev{\op{x}^{2}}{n}}$ et
  $\Delta p=\sqrt{\ev{\op{p}^{2}}{n}}$, d'où
  \begin{equation}
    \Delta x\,\Delta p=\frac{\hbar}{2}\,(2n+1)=\hbar\Bigl(n+\frac12\Bigr).
    \label{eq:DxDp}
  \end{equation}
  D'autre part $\ev{\comm{\op{x}}{\op{p}}}{n}=\ii\hbar$, donc le membre de
  droite de l'\cref{eq:robertson} vaut $\hbar/2$. Puisque $n+\tfrac12\geq\tfrac12$,
  l'inégalité est bien vérifiée pour tout $n$.
\end{solution}

\begin{corollaire}\label{cor:saturation}
  L'état fondamental $\ket{0}$ sature l'inégalité de Robertson pour $(\op{x},\op{p})$ :
  $\Delta x\,\Delta p=\hbar/2$.
\end{corollaire}

\begin{remarque}[Énergie du point zéro]\label{rem:zero}
  D'après l'\cref{eq:En}, $E_{0}=\hbar\omega/2\neq 0$ : la particule ne peut pas
  être au repos au fond du puits. Physiquement, annuler à la fois la
  position et l'impulsion violerait l'\cref{eq:DxDp}, qui impose
  $\Delta x\,\Delta p\geq\hbar/2$. L'énergie du point zéro est la
  conséquence directe du principe d'incertitude~\cite{heisenberg1927,robertson1929}%
  \footnote{Cette énergie ne s'annule pas à température nulle : c'est une propriété de
  l'état fondamental, et non un effet thermique.}
  (voir aussi le traitement de référence de~\cite{griffiths2018,cohen1973}).
\end{remarque}

%% ----- Question 4 -----
\begin{question}[3 pts]\label{q:virial}
  Montrer que $\ev{\op{p}^{2}/2m}{n}=\ev{\tfrac12 m\omega^{2}\op{x}^{2}}{n}=E_{n}/2$.
\end{question}

\begin{solution}
  Par l'\cref{eq:p2}, $\ev{\op{p}^{2}/2m}{n}=\tfrac{\hbar\omega}{4}(2n+1)$ ; par
  l'\cref{eq:x2}, $\ev{\tfrac12 m\omega^{2}\op{x}^{2}}{n}=\tfrac{\hbar\omega}{4}(2n+1)$.
  Ces deux valeurs coïncident et valent $E_{n}/2$ d'après l'\cref{eq:En} :
  c'est le théorème du viriel pour un potentiel quadratique.
\end{solution}

%% ----- Question 5 -----
\begin{question}[2 pts]\label{q:numerique}
  Un oscillateur a pour fréquence $\nu=\qty{1}{\tera\hertz}$. Calculer
  l'énergie du point zéro en $\unit{\milli\electronvolt}$ et le nombre d'onde
  correspondant $\tilde\nu=\nu/c$ en $\unit{\wavenumber}$.
\end{question}

\begin{solution}
  Avec $\omega=2\pi\nu$, on a $E_{0}=\hbar\omega/2=h\nu/2$. Pour
  $h=\qty{6.62607015e-34}{\joule\second}$ (valeur exacte) et
  $\qty{1}{\electronvolt}=\qty{1.602176634e-19}{\joule}$,
  \begin{equation}
    E_{0}=\frac{h\nu}{2}
    =\frac{\qty{6.62607015e-34}{\joule\second}\times\qty{1e12}{\hertz}}{2}
    =\qty{3.313e-22}{\joule}\approx\qty{2.07}{\milli\electronvolt}.
  \end{equation}
  Le nombre d'onde est
  $\tilde\nu=\nu/c=\qty{1e12}{\hertz}/\qty{2.99792458e10}{\centi\metre\per\second}
  \approx\qty{33.4}{\wavenumber}$.
\end{solution}
